\documentclass[11pt]{article}

\usepackage[margin=1in]{geometry}
\usepackage[T1]{fontenc}
\usepackage{lmodern}
\usepackage{hyperref}
\usepackage{enumitem}

\hypersetup{
  colorlinks=true,
  linkcolor=black,
  urlcolor=blue
}

\setlist[itemize]{leftmargin=*}

\title{IFT GPU Workstation Memo}
\author{CQMP}
\date{\today}

\begin{document}
\maketitle

\section{Hardware Configuration}

The workstations are Dell Pro Max Tower T2 systems installed from the shared
workstation image.  The currently installed reference machine, Matsubara, has
the following configuration:

\begin{itemize}
  \item CPU: Intel Core Ultra 7 265, 20 cores, one socket.
  \item Memory: 30 GiB RAM plus 8 GiB compressed swap.
  \item Local storage: one Kioxia XG10d 1 TB NVMe device.  The standard
        installation uses a 512 MB EFI partition, a 1 GB \texttt{/boot}
        partition, a 50 GB operating-system partition, and the remainder as
        \texttt{/data}.
  \item Graphics: integrated Intel Arrow Lake-S graphics plus an NVIDIA
        RTX 2000 Ada Generation GPU.
  \item Network: Intel I219-LM Ethernet.
\end{itemize}

The workstation inventory and assigned addresses are tracked in
\texttt{machines.md} in the image repository.

\section{System Configuration}

The machines run CentOS Stream 9 as a bootc image.  In this model, the
operating system is built as a container image and deployed atomically to the
host.  The machine normally boots the current deployment and keeps a rollback
deployment available.  Updates are delivered by publishing a new image and then
running a bootc upgrade on the workstation.

The image repository is:

\begin{center}
\url{https://github.com/CQMP/workstation-image}
\end{center}

The published image used by the installer and by bootc is:

\begin{center}
\texttt{ghcr.io/cqmp/centos9-workstation:latest}
\end{center}

The main image definition is the \texttt{Containerfile}.  The unattended
installer is described by \texttt{ks.cfg}; it partitions the local NVMe disk,
pulls the image above, and installs the bootc system.  The image includes GNOME
and GDM, NVIDIA drivers and CUDA tooling, SSH, LDAP/SSSD authentication,
autofs/NFS integration, HTCondor, compilers, MPI stacks, Python scientific
packages, browsers, editors, and common desktop tools.

The image also carries kernel arguments and initramfs configuration needed for
these systems.  In particular, nouveau is blacklisted, NVIDIA DRM mode setting
is enabled, and the initramfs is kept small enough for reliable EFI boot.

\section{Package Requests and Image Changes}

Users should not install long-lived system packages manually on individual
workstations.  Manual changes are local to one machine and may disappear after
reinstallation or diverge from the rest of the pool.  Instead, package requests
and system configuration changes should be made through the image repository.

The preferred workflow is:

\begin{enumerate}
  \item Fork \url{https://github.com/CQMP/workstation-image}.
  \item Modify the \texttt{Containerfile} or the relevant file under
        \texttt{etc/}, \texttt{usr/}, or another tracked directory.
  \item Build or inspect the change locally when practical.
  \item Submit a pull request against the main repository.
\end{enumerate}

After a pull request is merged, GitHub Actions rebuilds and publishes the image.
Machines can then receive the update with:

\begin{verbatim}
sudo bootc upgrade
sudo reboot
\end{verbatim}

\section{Logging In}

Authentication is handled through the FUW/ITP LDAP system via SSSD.  Users log
in with their Unix password.  Being present in LDAP is necessary but not
sufficient: the image also contains an explicit allow list of users who may log
in to these workstations.

The allow list is configured in \texttt{etc/sssd/sssd.conf} as
\texttt{simple\_allow\_users}.  Users who need access should provide their Unix
user name so it can be added to the image.  The current image allow list is:

\begin{center}
\texttt{egull, rfarid, ajazdzewska, amarie, abalbi}
\end{center}

Home directories and shared file systems are resolved through LDAP/autofs.  If
login fails, the first things to check are whether the user exists in the LDAP
Unix domain, whether the user is on the image allow list, and whether the user
is using the Unix password rather than another institutional password.

\section{Condor}

Each workstation is configured as an HTCondor execute node.  The local Condor
configuration is in:

\begin{center}
\texttt{etc/condor/config.d/00-ift-execute.conf}
\end{center}

The execute node contacts \texttt{condor.gull-group.org}.  It runs only the
Condor master and startd daemons, so these machines provide execution slots but
are not the central manager.

The configuration is designed for interactive workstations.  Jobs start only
when the keyboard has been idle for at least 15 minutes and the machine load is
low:

\begin{center}
\texttt{START = (KeyboardIdle > 900) \&\& (LoadAvg < 0.5)}
\end{center}

If a local user returns briefly, Condor suspends the job.  If the machine stays
active, the job is preempted and evicted.  The image enables CRIU-based
checkpointing and installs NVIDIA's \texttt{cuda-checkpoint} helper so GPU jobs
can be checkpointed during eviction when supported by the workload.  The
configured vacate time is one hour, giving checkpointing time to complete before
the job is killed.

The NVIDIA GPU is intended to be available to CUDA and Condor jobs.  Users
submitting GPU jobs should request the appropriate GPU resources through the
central Condor configuration and should test their jobs with \texttt{nvidia-smi}
or CUDA device discovery before launching large runs.

\section{File System}

Authentication and network file-system access are handled through ITP/FUW Unix
credentials.  The image configures SSSD for LDAP and autofs for the shared
mount points:

\begin{center}
\texttt{/dmj}, \texttt{/expo}, and \texttt{/repo}
\end{center}

These network file systems are appropriate for home directories, shared code,
small inputs, and files that need central management.

Each workstation also has a fast local NVMe scratch disk mounted at
\texttt{/data}.  The standard installation creates per-user directories there
for allowed users.  This disk is not backed up.  It should be used for
simulation scratch data, compilation trees, temporary files, local virtual
environments, and other high-I/O work that does not need backup.  Important
results should be copied back to backed-up network storage.

\section{Remote Desktop Access (VNC)}

The workstations run TigerVNC server.  Each user maintains a private VNC session
on a chosen display number (e.g., \texttt{:1} on port 5901).  Sessions persist
across disconnects: GUI applications keep running after the viewer is closed.

If two users share a machine they must use different display numbers
(\texttt{:1}, \texttt{:2}, \ldots); the corresponding TCP port is
$5900 + \text{display number}$.

\subsection*{First-time setup (on the workstation, via SSH)}

Set a VNC password:
\begin{verbatim}
vncpasswd
\end{verbatim}

Allow your user services to run at boot even when you are not logged in:
\begin{verbatim}
loginctl enable-linger
\end{verbatim}

Enable and start your VNC session:
\begin{verbatim}
systemctl --user enable --now vncserver@:1.service
\end{verbatim}

The session will now start automatically after every reboot.

\subsection*{Connecting from macOS}

VNC traffic is unencrypted; always route it through an SSH tunnel.  In a
terminal, forward the VNC port to your local machine:

\begin{verbatim}
ssh -N -L 5901:Matsubara:5901 <username>@<machine-ip>
\end{verbatim}

Keep that terminal open (or background it with \texttt{\&}).  Then open the
built-in macOS Screen Sharing app---press \textbf{Cmd+K} in Finder, or
search for ``Screen Sharing'' in Spotlight---and connect to:

\begin{verbatim}
vnc://localhost:5901
\end{verbatim}

\href{https://www.realvnc.com/en/connect/download/viewer/}{RealVNC Viewer}
(free download) is a more comfortable alternative: it offers better compression,
a saved-connections list, and clipboard integration.  After the SSH tunnel is
running, create a new connection to \texttt{localhost:5901}.

\subsection*{Session management}

\begin{verbatim}
systemctl --user stop    vncserver@:1.service   # stop the session
systemctl --user disable vncserver@:1.service   # prevent autostart on reboot
\end{verbatim}

\section{Matrix/Element Messaging Server --- Installation Status (June 2026)}

\subsection*{Architecture}

A self-hosted Matrix homeserver is being set up to replace the group's Slack
workspace (free plan, message history limited to 90 days).  The deployment
consists of two components:

\begin{itemize}
  \item \textbf{Tuwunel v1.7.1} (Rust Matrix homeserver) running on an Oracle
        Cloud Always-Free VM (\texttt{gull-matrix-vnic}, 1 OCPU / 1 GB RAM,
        Frankfurt region, public IP \texttt{79.76.114.255}).  DNS A record
        \texttt{matrix.gull-group.org} points to this machine.  Matrix user
        IDs will take the form \texttt{@user:matrix.gull-group.org}.

  \item \textbf{MinIO} (S3-compatible object store) running in Docker on the
        Synology NAS at \texttt{fermi.physics.lsa.umich.edu} (DS2419+, DSM
        7.3.2), providing persistent media storage backed by local NAS disks.
\end{itemize}

The Synology sits behind the University of Michigan campus firewall (inbound
connections blocked).  A persistent reverse SSH tunnel --- maintained by an
\texttt{autossh} container on the Synology --- forwards
\texttt{127.0.0.1:9000} on the Oracle VM through to MinIO on the Synology.
Tuwunel reads and writes all uploaded media (files, images, avatars) via this
tunnel.  MinIO is never exposed directly to the internet.

\subsection*{Completed}

\begin{itemize}
  \item Oracle VM provisioned; 2 GB swap file added; OS firewall opened on
        TCP 80, 443 and UDP 51820.
  \item DNS A record \texttt{matrix.gull-group.org $\to$ 79.76.114.255} created.
  \item Restricted \texttt{tunnel} SSH user created on Oracle VM; autossh
        container on Synology holds the reverse tunnel open
        (\texttt{PermitOpen 127.0.0.1:9000} enforced in \texttt{sshd\_config}).
  \item MinIO running on Synology (Docker, host networking, bound to
        \texttt{127.0.0.1:9000}); bucket \texttt{matrix-media} created;
        MinIO credentials stored in
        \texttt{/volume1/docker/matrix/docker-compose.yml} on the Synology.
  \item Tuwunel v1.7.1 installed via \texttt{.deb} on Oracle VM; systemd
        service enabled.  Configured with \texttt{server\_name =
        "matrix.gull-group.org"}, \texttt{allow\_registration = false}, and
        S3 media storage pointing at MinIO via the tunnel.
\end{itemize}

\subsection*{Remaining Tasks}

\begin{enumerate}
  \item \textbf{Fix S3 initialisation error.}  Tuwunel starts and listens on
        \texttt{127.0.0.1:8008} but logs a builder error when connecting to
        the S3 provider at startup.  Connectivity from the Oracle VM to MinIO
        is confirmed (\texttt{curl http://127.0.0.1:9000/minio/health/live}
        succeeds); the issue is in the S3 config block in
        \texttt{/etc/tuwunel/tuwunel.toml}.

  \item \textbf{Nginx reverse proxy.}  Install Nginx on the Oracle VM; proxy
        \texttt{https://matrix.gull-group.org} and the Matrix federation port
        to Tuwunel on \texttt{127.0.0.1:8008}.

  \item \textbf{TLS certificate.}  Obtain a Let's Encrypt certificate via
        Certbot for \texttt{matrix.gull-group.org}.

  \item \textbf{Federation check.}  Verify federation works using the Matrix
        federation tester (\url{https://federationtester.matrix.org}).

  \item \textbf{Create admin account.}  Register the first user
        (\texttt{egull}) via the Tuwunel admin API; user self-registration is
        disabled.

  \item \textbf{Nightly rsync backup.}  Add a cron job on the Oracle VM to
        rsync \texttt{/var/lib/tuwunel} to the Synology over the SSH tunnel
        as a safety net for the local RocksDB database.

  \item \textbf{User migration.}  Invite group members; communicate login
        instructions and recommended client (Element desktop, installed on all
        workstations).
\end{enumerate}

\subsection*{Key Locations}

\begin{center}
\begin{tabular}{ll}
\hline
Oracle VM SSH        & \texttt{ubuntu@79.76.114.255} \\
Tuwunel config       & \texttt{/etc/tuwunel/tuwunel.toml} on Oracle VM \\
Tuwunel data         & \texttt{/var/lib/tuwunel} on Oracle VM \\
Synology compose     & \texttt{/volume1/docker/matrix/docker-compose.yml} \\
MinIO data           & \texttt{/volume1/minio/data} on Synology \\
Tunnel key (Synology)& \texttt{/volume1/docker/matrix-tunnel/tunnel\_key} \\
\hline
\end{tabular}
\end{center}

\end{document}
